\subsection{Techniques}

\begin{itemize}
    \item \textbf{Check Even/Odd}: \texttt{$n \& 1$} can be used to check if a number is even (result is 0) or odd (result is 1).
    \item \textbf{Power of 2 check: } \texttt{$(n \&\& !(n \& (n - 1)))$} returns true if $n$ is a power of 2.
    \item \textbf{Count 1s: } \texttt{\_\_builtin\_popcount(n)} counts set bits in O(1). Use \texttt{\_\_builtin\_popcountll(n)} for \texttt{long long}. Manually: loop \texttt{$n \& (n-1)$} until $n = 0$, each iteration removes the rightmost set bit.
    \item \textbf{Isolate/Clear Rightmost Set Bit: } Isolate: \texttt{$n \& (-n)$}. Clear: \texttt{$n \& (n - 1)$}.
    \item \textbf{GCC Built-ins (O(1)):} \texttt{\_\_builtin\_clz(n)} — count leading zeros; \texttt{\_\_builtin\_ctz(n)} — count trailing zeros; \texttt{\_\_builtin\_parity(n)} — parity of set bits (0 or 1); \texttt{\_\_lg(n)} = $\lfloor \log_2 n \rfloor$. All undefined for $n = 0$.
    \item \textbf{Subset Generation: } Loop from $0$ to $2^n - 1$; each integer's bits represent inclusion/exclusion. To iterate over all subsets of a bitmask $m$: \texttt{for(int s=m; s>0; s=(s-1)\&m)}.
    \item \textbf{Lowest Set Bit index: } \texttt{\_\_builtin\_ctz(n)} gives the index of the lowest set bit directly.
\end{itemize}